% !TEX root = ../main.tex
Pennysia is an open liquidity network, a single exchange layer where every source of liquidity, wherever it sits, competes for the user's order. Decentralized exchanges, aggregators, intent solvers, market makers, perpetual and prediction protocols, vaults, fiat ramps, and licensed venues all answer the same auction. The best executable price wins. The user signs once. The network routes, executes, and verifies the fill. For builders, a venue connects once and competes for every order from day one.

\noindent Pennysia aims to combine the performance, liquidity access, and usability of centralized exchanges (CEXs) with the self-custody of decentralized exchanges (DEXs). The complexity of wallet management, gas payments, and chain routing is abstracted away, while users retain their choice of wallet and chain. Fiat enters and exits on card and bank rails. Liquidity is organized into two lanes: a \textbf{permissionless lane} of open venues, onchain and offchain, open to everyone, and a \textbf{regulated lane} of gated liquidity (centralized market makers, licensed venues, compliant pools, fiat ramps) unlocked by optional verification.

\noindent An experimental deployment on Ethereum proved the model: 30 venues competing in parallel per swap, atomic in-transaction fallback, and a settlement layer that enforces signed terms on every order. Pennysia v1, the Open Liquidity Network, generalizes that proof to every chain and every service class, a single interface for the world's liquidity.
